% Lecture notes: Week 2, Lecture 3 -- Stress
% To hide this lecture from the website, add a line containing only:  % draft
% To set the date shown on the website, add a line like:  % posted: 2026-09-02
\documentclass[11pt]{article}
\usepackage{mech2030notes}
\begin{document}

% ##################################################################
%  WEEK 2, LECTURE 3
% ##################################################################
\lecture{2}{3}{Stress}

\begin{center}
\begin{tikzpicture}[baseline=(b)]
  \coordinate (b) at (0,0);
  \dogbone{0}{0.3}{1.0}
  \node[pin] at (0,1.9) {}; \node[right] at (0.05,1.9) {$C$};
  \draw[thick, fill=white, fill opacity=0.6] (-0.75,-0.45) -- (0.45,-0.45) -- (0.75,0.05) -- (-0.45,0.05) -- cycle;
  \node[pin] at (0,-0.2) {}; \node[right] at (0.35,-0.2) {$B$};
  \draw[force] (0,2.55) -- (0,3.4) node[above] {$P$};
  \draw[force] (0,-2.55) -- (0,-3.4) node[below] {$P$};
  \draw[thinarrow] (2.6,0.1) -- (0.85,-0.15);
  \node[above] at (2.6,0.1) {\small cut plane};
\end{tikzpicture}
\hspace{2.5cm}
\begin{tikzpicture}[baseline=(b)]
  \coordinate (b) at (0,0);
  \draw[thick, fill=gray!8]
     (-0.6,0.6) -- (-0.6,-1.0) to[out=-90,in=90] (-1.1,-1.6) -- (-1.1,-2.5)
     -- (1.1,-2.5) -- (1.1,-1.6) to[out=90,in=-90] (0.6,-1.0) -- (0.6,0.6);
  \draw[thick, fill=gray!20] (0,0.6) ellipse (0.6 and 0.18);
  \foreach \x in {-0.45,-0.22,0,0.22,0.45}
    \draw[-{Stealth[length=2mm,width=1.5mm]}, semithick] (\x,0.6) -- (\x,1.4);
  \node[right] at (0.55,1.3) {$\sigma$};
  \draw[force] (0,1.7) -- (0,2.6) node[above] {$N$};
  \draw[force] (0,-2.55) -- (0,-3.4) node[below] {$P$};
  \draw[thinarrow] (2.2,0.2) -- (0.65,0.55);
  \node[right, align=left, font=\small] at (2.2,0.2) {cross-section\\area $A$};
\end{tikzpicture}
\end{center}

\begin{itemize}[label=$\rightarrow$]
  \item Generally, $\displaystyle N = \int_A \sigma\, dA$.
  \item If $\sigma$ is constant, $N = \sigma A$. Thus
        \[ \boxed{\sigma = \frac{N}{A}} \quad \text{is the \emph{average normal stress}.} \]
  \item For the dogbone, $\sigma_B = \dfrac{P}{A_B}$ and $\sigma_C = \dfrac{P}{A_C}$.
  \begin{itemize}[label=\then]
    \item Since $A_B < A_C$, $\;\Rightarrow\; \sigma_B > \sigma_C$.
    \item Fracture occurs at $B$, in the thinner section.
  \end{itemize}
\end{itemize}

\subsection*{Complete (2D) internal forces}

\noindent
\begin{minipage}[c]{0.45\textwidth}
\centering
\begin{tikzpicture}
  \lwall{0}{-0.8}{0.8}
  \cutbar{0}{2.4}{0.2}{2}
  \draw[force] (2.75,0.4) -- (2.75,-0.6) node[below] {$V$};
  \draw[force] (2.6,0) -- (3.5,0) node[right] {$N$};
  \ccwm{4.2}{0}{$M$}
\end{tikzpicture}
\end{minipage}%
\hfill
\begin{minipage}[c]{0.52\textwidth}
\begin{itemize}
  \item $N$ causes a \emph{normal} stress: $\;\boxed{\sigma = N/A}$
  \item $V$ causes a \emph{shear} stress: $\;\boxed{\tau = V/A}$
  \item $M$ causes a \emph{distribution} of normal stresses (more on this later).
\end{itemize}
\end{minipage}

\section*{Example: Two cylinders}

\noindent
\begin{minipage}[c]{0.62\textwidth}
\centering
\begin{tikzpicture}
  \draw[beam] (0,-0.2) rectangle (2.4,0.2);
  \draw[beam] (2.4,-0.8) rectangle (5,0.8);
  \node[pin] at (1.2,0) {}; \node[above] at (1.2,0.55) {$A$};
  \node[pin] at (3.9,0) {}; \node[above right] at (3.9,0) {$B$};
  \draw[force] (-1.3,0) node[left] {$\SI{50}{N}$} -- (-0.05,0);
  \draw[force] (3.6,0) -- (2.5,0);
  \node[above] at (3.05,0.05) {\small$\SI{150}{N}$};
  \draw[force] (5.05,0) -- (6.2,0) node[right] {$\SI{100}{N}$};
  \draw[thinarrow] (1.6,0.75) -- (1.6,0.22);
  \draw[thinarrow] (1.6,-0.75) -- (1.6,-0.22);
  \node[below] at (1.6,-0.75) {\small$\SI{3}{mm}$};
  \draw[thinarrow] (4.5,1.35) -- (4.5,0.82);
  \draw[thinarrow] (4.5,-1.35) -- (4.5,-0.82);
  \node[below] at (4.5,-1.35) {\small$\SI{10}{mm}$};
\end{tikzpicture}
\end{minipage}%
\hfill
\begin{minipage}[c]{0.34\textwidth}
\textbf{Q:} What are the stresses at $A$ and $B$?
\end{minipage}

\subsection*{``Cut'' at $A$}

\noindent
\begin{minipage}[c]{0.45\textwidth}
\centering
\begin{tikzpicture}
  \cutbar{0}{1.4}{0.2}{2}
  \draw[force] (-1.3,0) node[left] {$\SI{50}{N}$} -- (-0.05,0);
  \draw[force] (1.5,0) -- (2.5,0) node[right] {$N_A$};
\end{tikzpicture}
\end{minipage}%
\hfill
\begin{minipage}[c]{0.52\textwidth}
\[
  \textstyle\sum F_x = 0 \;\rightarrow\; N_A + \SI{50}{N} = 0, \qquad N_A = -\SI{50}{N}
\]
\end{minipage}

\[
  \boxed{\sigma_A = \frac{N_A}{A_A} = \frac{-\SI{50}{N}}{\pi(\SI{1.5}{mm})^2}
         = -\SI{7.1}{MPa}} \quad \text{(compression)}
\]

\paragraph{Note on units:}
\[
  \si{Pa} = \frac{\si{N}}{\si{m^2}}, \qquad
  \frac{\si{N}}{\si{mm^2}}\left(\frac{\SI{1e3}{mm}}{\SI{1}{m}}\right)^2
  = \num{1e6}\,\frac{\si{N}}{\si{m^2}} = \si{MPa}
\]
\begin{itemize}[label=\then]
  \item Working in \si{N} and \si{mm} gives stresses directly in \si{MPa}.
\end{itemize}

\subsection*{``Cut'' at $B$}

\noindent
\begin{minipage}[c]{0.45\textwidth}
\centering
\begin{tikzpicture}
  \cutbar{0}{1.6}{0.8}{1}
  \draw[force] (-0.1,0) -- (-1.1,0) node[left] {$N_B$};
  \draw[force] (1.65,0) -- (2.6,0) node[right] {$\SI{100}{N}$};
\end{tikzpicture}
\end{minipage}%
\hfill
\begin{minipage}[c]{0.52\textwidth}
\[
  \textstyle\sum F_x = 0 \;\rightarrow\; -N_B + \SI{100}{N} = 0, \qquad N_B = \SI{100}{N}
\]
\end{minipage}

\[
  \boxed{\sigma_B = \frac{N_B}{A_B} = \frac{\SI{100}{N}}{\pi(\SI{5}{mm})^2}
         = \SI{1.27}{MPa}} \quad \text{(tension)}
\]

\end{document}
